\paper{非数学 A 类}

\begin{problems}
\item[一、]（本题 30 分，每小题 6 分）填空题：
\begin{enumerate}
  \item 设函数 $f(x)=x-\ln(1+x)-\frac{1}{2}x\sin x$，$g(x)=ax^3$，且满足 $\displaystyle \lim_{x\to 0}\frac{f(x)}{g(x)}=1$，则 $a=\underline{\hspace{3.5em}}$．
  \item 设平面过原点与点 $(6,-3,2)$，且与平面 $4x-y+2z-8=0$ 垂直，则此平面方程为 $\underline{\hspace{3.5em}}$．
  \item 极限 $\displaystyle \lim_{x\to 0}\frac{\displaystyle \int_0^{\sin^2 x}\ln(1+t)\dd t}{(\sqrt[3]{1+x^3}-1)\sin x}=\underline{\hspace{3.5em}}$．
  \item 设函数 $f(u,v)$ 具有一阶连续偏导数，$f(2,0)=3$，$f_u'(2,0)=5$，又设 $z=z(x,y)$ 是由方程 $xz=f(2x-y,xyz)$ 确定的隐函数，则 $\displaystyle \left.\frac{\partial z}{\partial x}\right|_{\substack{x=1\\ y=0}}=\underline{\hspace{3.5em}}$．
  \item 设曲线 $L$ 为球面 $x^2+y^2+z^2=4$ 与平面 $x+y+z=3$ 的交线，则曲线积分
  \[
    \int_L (x^2+2xy)\dd s=\underline{\hspace{10em}}.
  \]
\end{enumerate}

\item[二、]（本题 14 分）设 $f(x)$ 在 $(-\infty,+\infty)$ 上连续，在 $x=0$ 处可导，且满足
\[
f(x)=x^3+x^2\lim_{x\to 0}\frac{f(x)}{x}-x\int_0^1 f(x)\dd x.
\]
求函数 $f(x)$ 的表达式．

\item[三、]（本题 14 分）求微分方程 $y'x\ln x\sin y+(1-x\cos y)\cos y=0$ 的通解．

\item[四、]（本题 14 分）设 $\Omega$ 是由平面 $\Pi_1:\ x=0$，$\Pi_2:\ x=y$，$\Pi_3:\ y=z$，$\Pi_4:\ z=1$，$\Pi_5:\ z=2$ 围成的空间区域，求三重积分
\[
\iiint_{\Omega}\frac{\dd x\dd y\dd z}{\sqrt{y^2+z^2}}.
\]

\item[五、]（本题 14 分）判断级数 $\displaystyle \sum_{n=1}^{\infty}a_n$ 的收敛性，其中 $a_n=\sin^2 x\sin^2 2x\cdots\sin^2 2^n x$，\mbox{$x\in(-\infty,+\infty)$}．

\item[六、]（本题 14 分）设函数 $f(x),g(x)$ 在 $[a,b]$ 上连续可导且满足以下三个条件：
\begin{enumerate}[label=(\roman*)]
  \item $\displaystyle \int_a^b f(x)\dd x=\int_a^b g(x)\dd x=0$；
  \item $f'(x)g(x)-f(x)g'(x)>0$，$x\in[a,b]$；
  \item $g(x)$ 在 $(a,b)$ 上有唯一零点 $\xi$．
\end{enumerate}
证明：
\begin{enumerate}
  \item $\displaystyle \int_a^b f(x)\left(\int_a^x g(s)\dd s\right)\dd x=-\int_a^b g(x)\left(\int_a^x f(s)\dd s\right)\dd x$．
  \item $\displaystyle \int_a^b f(x)\left(\int_a^x g(s)\dd s\right)\dd x>0$．
\end{enumerate}
\end{problems}
